\documentclass{ximera}

\input{../../../preamble.tex}


\definecolor{fvdnavy}{HTML}{071D43}
\definecolor{fvdcyan}{HTML}{20BFC6}
\definecolor{fvdcyanlight}{HTML}{EFFCFB}
\definecolor{fvdmagenta}{HTML}{F10D69}
\definecolor{fvdmagentalight}{HTML}{FFF1F7}
\definecolor{fvdtext}{HTML}{163044}





%=================================================
% Pedagogische omgevingen
% PDF: tcolorbox
% HTML: learning-box
%=================================================

\newenvironment{voorkennis}
{%
    \ifdefined\HCode
        \HCode{%
<section class="learning-box learning-box--prior">
<div class="learning-box__header">
<span class="learning-box__icon" aria-hidden="true"></span>
<span class="learning-box__title">Voorkennis</span>
</div>
<div class="learning-box__content">
        }%
        \begin{itemize}
    \else
       \begin{tcolorbox}[
    enhanced,
    breakable,
    colback=fvdcyanlight,
    colframe=fvdcyan,
    coltext=fvdtext,
    boxrule=0.5pt,
    leftrule=4pt,
    arc=4mm,
    outer arc=4mm,
    left=7mm,
    right=6mm,
    top=5mm,
    bottom=5mm,
    before skip=6mm,
    after skip=6mm,
    title=Voorkennis,
    coltitle=fvdcyan,
    fonttitle=\bfseries\Large,
    attach title to upper={\par\smallskip},
    boxed title style={
        empty,
        left=0mm,
        right=0mm,
        top=0mm,
        bottom=0mm
    },
    overlay={
        \draw[
            fvdcyan,
            line width=0.5pt
        ]
        ([xshift=42mm,yshift=-13mm]frame.north west)
        --
        ([xshift=-7mm,yshift=-13mm]frame.north east);
    }
]
        \begin{itemize}
    \fi
}
{%
    \end{itemize}
    \ifdefined\HCode
        \HCode{%
</div>
</section>
        }%
    \else
        \end{tcolorbox}
    \fi
}


\newenvironment{leerdoelen}
{%
    \ifdefined\HCode
        \HCode{%
<section class="learning-box learning-box--goals">
<div class="learning-box__header">
<span class="learning-box__icon" aria-hidden="true"></span>
<span class="learning-box__title">Leerdoelen</span>
</div>
<div class="learning-box__content">
        }%
        \begin{itemize}
    \else
\begin{tcolorbox}[
    enhanced,
    breakable,
    colback=fvdmagentalight,
    colframe=fvdmagenta,
    coltext=fvdtext,
    boxrule=0.5pt,
    leftrule=4pt,
    arc=4mm,
    outer arc=4mm,
    left=7mm,
    right=6mm,
    top=5mm,
    bottom=5mm,
    before skip=6mm,
    after skip=6mm,
    title=Leerdoelen,
    coltitle=fvdmagenta,
    fonttitle=\bfseries\Large,
    attach title to upper={\par\smallskip},
    boxed title style={
        empty,
        left=0mm,
        right=0mm,
        top=0mm,
        bottom=0mm
    },
    overlay={
        \draw[
            fvdmagenta,
            line width=0.5pt
        ]
        ([xshift=42mm,yshift=-13mm]frame.north west)
        --
        ([xshift=-7mm,yshift=-13mm]frame.north east);
    }
]
        \begin{itemize}
    \fi
}
{%
    \end{itemize}
    \ifdefined\HCode
        \HCode{%
</div>
</section>
        }%
    \else
        \end{tcolorbox}
    \fi
}


\addPrintStyle{../../..}

\title{Oefeningen — De kettingregel}
\author{Ingmar Herreman}

\begin{document}

% FVD_CHAPTER_DATA_START
% Automatisch gegenereerd uit index.tex.
% Niet handmatig aanpassen.
\fvdchapterdata{3}{Afleiden van samengestelde functie}{2}
% FVD_CHAPTER_DATA_END




\maketitle



\begin{oefening}[Binnen en buiten herkennen]{\oefB\ESS}

Bepaal telkens de binnenste en de buitenste functie.

\begin{enumerate}
  \item \[f(x)=(3x-2)^7.\]
  \item \[g(x)=\sqrt{x^2+5x+8}.\]
  \item \[h(x)=\frac{1}{(2x+1)^4}.\]
  \item \[p(x)=\left(x^3-4x+1\right)^5.\]
  \item \[q(x)=\sqrt{(x-1)^2+4}.\]
\end{enumerate}

\begin{oplossing}
\begin{enumerate}
  \item Binnen: \[u=3x-2.\] Buiten: \[f(u)=u^7.\]
  \item Binnen: \[u=x^2+5x+8.\] Buiten: \[g(u)=\sqrt{u}.\]
  \item Binnen: \[u=2x+1.\] Buiten: \[h(u)=u^{-4}.\]
  \item Binnen: \[u=x^3-4x+1.\] Buiten: \[p(u)=u^5.\]
  \item Buiten: \[\sqrt{\phantom{x}}.\] Binnen: \[(x-1)^2+4.\]
  De binnenste uitdrukking bevat zelf opnieuw een samenstelling:
  \[(x-1)^2.\]
\end{enumerate}
\end{oplossing}

\end{oefening}



\begin{oefening}[De kettingregel toepassen]{\oefB\ESS}

Bepaal telkens de afgeleide.

\begin{enumerate}
  \item \[f(x)=(5x+2)^4.\]
  \item \[g(x)=(x^2-3)^6.\]
  \item \[h(x)=\sqrt{4x+7}.\]
  \item \[p(x)=(2x^2+x-1)^5.\]
  \item \[q(x)=\frac{1}{(3x-2)^3}.\]
  \item \[r(x)=\sqrt{x^3+1}.\]
\end{enumerate}

\begin{oplossing}
\begin{enumerate}
  \item \[f'(x)=4(5x+2)^3\cdot5.\]
  Dus:
  \[\boxed{f'(x)=20(5x+2)^3.}\]

  \item \[g'(x)=6(x^2-3)^5\cdot2x.\]
  Dus:
  \[\boxed{g'(x)=12x(x^2-3)^5.}\]

  \item \[h'(x)=\frac{1}{2\sqrt{4x+7}}\cdot4.\]
  Dus:
  \[\boxed{h'(x)=\frac{2}{\sqrt{4x+7}}.}\]

  \item \[p'(x)=5(2x^2+x-1)^4(4x+1).\]

  \item We schrijven:
  \[q(x)=(3x-2)^{-3}.\]
  Dan:
  \[q'(x)=-3(3x-2)^{-4}\cdot3.\]
  Dus:
  \[\boxed{q'(x)=-\frac{9}{(3x-2)^4}.}\]

  \item \[r'(x)=\frac{1}{2\sqrt{x^3+1}}\cdot3x^2.\]
  Dus:
  \[\boxed{r'(x)=\frac{3x^2}{2\sqrt{x^3+1}}.}\]
\end{enumerate}
\end{oplossing}

\end{oefening}




\begin{oefening}[Meerdere lagen]{\oefB\ESS}

Bepaal de afgeleide.

\begin{enumerate}
  \item \[f(x)=\left((2x-1)^2+3\right)^4.\]
  \item \[g(x)=\sqrt{(x+2)^3+1}.\]
  \item \[h(x)=\left((x^2+1)^3-2\right)^5.\]
  \item \[p(x)=\frac{1}{\left((3x-1)^2+4\right)^2}.\]
\end{enumerate}

\begin{oplossing}
\begin{enumerate}
  \item
  \[f'(x)=4\left((2x-1)^2+3\right)^3\cdot2(2x-1)\cdot2.\]
  Dus:
  \[\boxed{f'(x)=16(2x-1)\left((2x-1)^2+3\right)^3.}\]

  \item
  \[g'(x)=\frac{1}{2\sqrt{(x+2)^3+1}}\cdot3(x+2)^2.\]
  Dus:
  \[\boxed{g'(x)=\frac{3(x+2)^2}{2\sqrt{(x+2)^3+1}}.}\]

  \item
  \[h'(x)=5\left((x^2+1)^3-2\right)^4\cdot3(x^2+1)^2\cdot2x.\]
  Dus:
  \[\boxed{h'(x)=30x(x^2+1)^2\left((x^2+1)^3-2\right)^4.}\]

  \item
  \[p(x)=\left((3x-1)^2+4\right)^{-2}.\]
  Dus:
  \[p'(x)=-2\left((3x-1)^2+4\right)^{-3}\cdot2(3x-1)\cdot3.\]
  Daarom:
  \[\boxed{p'(x)=-\frac{12(3x-1)}{\left((3x-1)^2+4\right)^3}.}\]
\end{enumerate}
\end{oplossing}

\end{oefening}



\begin{oefening}[Domein en afgeleide]{\oefB\ESS}

Bepaal voor elke functie:
\begin{enumerate}
  \item het domein van de functie;
  \item de afgeleide;
  \item het domein van de afgeleide.
\end{enumerate}

\begin{question}
\[f(x)=\sqrt{5-2x}.\]
\end{question}

\begin{question}
\[g(x)=\sqrt{x^2-4}.\]
\end{question}

\begin{question}
\[h(x)=\frac{1}{(x-3)^4}.\]
\end{question}

\begin{oplossing}

Voor
\[f(x)=\sqrt{5-2x}\]
moet:
\[5-2x\geq0.\]

Dus:
\[\operatorname{dom}f=\left]-\infty,\frac52\right].\]

Verder:
\[f'(x)=-\frac{1}{\sqrt{5-2x}}.\]

Daarom:
\[\operatorname{dom}f'=\left]-\infty,\frac52\right[.\]

Voor
\[g(x)=\sqrt{x^2-4}\]
moet:
\[x^2-4\geq0.\]

Dus:
\[\operatorname{dom}g=]-\infty,-2]\cup[2,+\infty[.\]

Verder:
\[g'(x)=\frac{x}{\sqrt{x^2-4}}.\]

Daarom:
\[\operatorname{dom}g'=]-\infty,-2[\cup]2,+\infty[.\]

Voor
\[h(x)=\frac{1}{(x-3)^4}\]
geldt:
\[\operatorname{dom}h=\mathbb R\setminus\{3\}.\]

We schrijven:
\[h(x)=(x-3)^{-4}.\]

Dus:
\[h'(x)=-4(x-3)^{-5}.\]

Daarom:
\[\boxed{h'(x)=-\frac4{(x-3)^5}}\]
met:
\[\operatorname{dom}h'=\mathbb R\setminus\{3\}.\]

\end{oplossing}

\end{oefening}



\begin{oefening}[Welke regel eerst?]{\oefU\ESS}

Bepaal eerst de hoofdstructuur van de functie.
Noteer welke afleidingsregel je eerst gebruikt en bereken daarna de afgeleide.

\begin{enumerate}
  \item \[f(x)=x^3(x^2+1)^4.\]
  \item \[g(x)=\frac{(2x+1)^5}{x^2+3}.\]
  \item \[h(x)=\sqrt{x^2+4x+7}.\]
  \item \[p(x)=\frac{x^2}{\sqrt{x+1}}.\]
  \item \[q(x)=(x^2-1)^3(x+2)^4.\]
\end{enumerate}

\begin{oplossing}
\begin{enumerate}
  \item Hoofdregel: productregel.
  \[
  f'(x)=3x^2(x^2+1)^4+x^3\cdot4(x^2+1)^3\cdot2x.
  \]
  Dus:
  \[\boxed{f'(x)=x^2(x^2+1)^3(11x^2+3).}\]

  \item Hoofdregel: quotiëntregel.
  \[
  g'(x)=
  \frac{10(2x+1)^4(x^2+3)-(2x+1)^5(2x)}{(x^2+3)^2}.
  \]
  Dus:
  \[\boxed{g'(x)=\frac{(2x+1)^4(6x^2-2x+30)}{(x^2+3)^2}.}\]

  \item Hoofdregel: kettingregel.
  \[
  h'(x)=\frac{2x+4}{2\sqrt{x^2+4x+7}}.
  \]
  Dus:
  \[\boxed{h'(x)=\frac{x+2}{\sqrt{x^2+4x+7}}.}\]

  \item Efficiënt is herschrijven:
  \[p(x)=x^2(x+1)^{-1/2}.\]
  Dan:
  \[
  p'(x)=2x(x+1)^{-1/2}-\frac12x^2(x+1)^{-3/2}.
  \]
  Dus:
  \[\boxed{p'(x)=\frac{x(3x+4)}{2(x+1)^{3/2}}.}\]

  \item Hoofdregel: productregel.
  \[
  q'(x)=6x(x^2-1)^2(x+2)^4+4(x^2-1)^3(x+2)^3.
  \]
  Dus:
  \[\boxed{q'(x)=2(x^2-1)^2(x+2)^3(5x^2+6x-2).}\]
\end{enumerate}
\end{oplossing}

\end{oefening}




\begin{oefening}[Wat ging er fout?]{\oefU\ESS}

Analyseer elke afleiding.
Geef aan of ze correct is. Indien niet, leg de fout uit en verbeter.

\begin{enumerate}
  \item
  \[
  f(x)=(3x+1)^4,\qquad
  f'(x)=4(3x+1)^3.
  \]

  \item
  \[
  g(x)=(x^2+1)^5,\qquad
  g'(x)=5(2x)^4\cdot2.
  \]

  \item
  \[
  h(x)=x(x+1)^3,\qquad
  h'(x)=3x(x+1)^2.
  \]

  \item
  \[
  p(x)=\sqrt{2x^2+3},\qquad
  p'(x)=\frac{4x}{2\sqrt{2x^2+3}}.
  \]
\end{enumerate}

\begin{oplossing}
\begin{enumerate}
  \item Fout. De afgeleide van de binnenste functie ontbreekt.
  \[
  f'(x)=4(3x+1)^3\cdot3.
  \]
  Dus:
  \[\boxed{f'(x)=12(3x+1)^3.}\]

  \item Fout. Bij het afleiden van de buitenste macht moet $x^2+1$ blijven staan.
  \[
  g'(x)=5(x^2+1)^4\cdot2x.
  \]
  Dus:
  \[\boxed{g'(x)=10x(x^2+1)^4.}\]

  \item Fout. De hoofdstructuur is een product.
  \[
  h'(x)=(x+1)^3+3x(x+1)^2.
  \]
  Dus:
  \[\boxed{h'(x)=(x+1)^2(4x+1).}\]

  \item Correct.
  \[
  p'(x)=\frac{1}{2\sqrt{2x^2+3}}\cdot4x
  =\frac{2x}{\sqrt{2x^2+3}}.
  \]
\end{enumerate}
\end{oplossing}

\end{oefening}




\begin{oefening}[Afleiden]{\oefU\ESS}

Bepaal de afgeleide.

\begin{enumerate}
  \item \[f(x)=\frac{\sqrt{x^2+1}}{x+2}.\]
  \item \[g(x)=\frac{(x^2-1)^3}{(2x+1)^2}.\]
  \item \[h(x)=x^2\sqrt{3x-1}.\]
  \item \[p(x)=\frac{\sqrt{(x-1)^2+4}}{(x+1)^3}.\]
\end{enumerate}

\begin{oplossing}
\begin{enumerate}
  \item
  \[
  f'(x)=
  \frac{\frac{x}{\sqrt{x^2+1}}(x+2)-\sqrt{x^2+1}}{(x+2)^2}.
  \]
  Dus:
  \[\boxed{f'(x)=\frac{2x-1}{(x+2)^2\sqrt{x^2+1}}.}\]

  \item
  \[
  g'(x)=
  \frac{6x(x^2-1)^2(2x+1)^2-(x^2-1)^3\cdot4(2x+1)}
       {(2x+1)^4}.
  \]
  Dus:
  \[\boxed{g'(x)=\frac{2(x^2-1)^2(x^2+3x+2)}{(2x+1)^3}.}\]

  \item
  \[
  h'(x)=2x\sqrt{3x-1}+x^2\frac{3}{2\sqrt{3x-1}}.
  \]
  Dus:
  \[\boxed{h'(x)=\frac{x(15x-4)}{2\sqrt{3x-1}}.}\]

  \item
  Met
  \[
  u(x)=\sqrt{(x-1)^2+4}
  \]
  geldt:
  \[
  u'(x)=\frac{x-1}{\sqrt{(x-1)^2+4}}.
  \]
  Met
  \[
  v(x)=(x+1)^3
  \]
  geldt:
  \[
  v'(x)=3(x+1)^2.
  \]
  Dus:
  \[
  p'(x)=
  \frac{
  \frac{x-1}{\sqrt{(x-1)^2+4}}(x+1)^3
  -
  3(x+1)^2\sqrt{(x-1)^2+4}
  }
  {(x+1)^6}.
  \]
\end{enumerate}
\end{oplossing}

\end{oefening}




\begin{oefening}[Een ballon die groter wordt]{\oefU\ESS}

Een bolvormige ballon heeft volume
\[
V=\frac{4}{3}\pi r^3.
\]

De straal verandert in de tijd volgens
\[
r(t)=2+0{,}4t,
\]
waarbij $r$ in centimeter en $t$ in seconden wordt uitgedrukt.

\begin{question}
Schrijf $V$ als functie van $t$.
\end{question}

\begin{question}
Bepaal $\dfrac{dV}{dt}$.
\end{question}

\begin{question}
Bereken hoe snel het volume verandert op het tijdstip $t=5$.
\end{question}

\begin{question}
Leg uit waar in deze berekening de kettingregel optreedt.
\end{question}

\begin{oplossing}

\[
V(t)=\frac{4}{3}\pi(2+0{,}4t)^3.
\]

Met de kettingregel:
\[
\frac{dV}{dt}
=
\frac{4}{3}\pi\cdot3(2+0{,}4t)^2\cdot0{,}4.
\]

Dus:
\[
\boxed{\frac{dV}{dt}=1{,}6\pi(2+0{,}4t)^2.}
\]

Voor $t=5$:
\[
r(5)=4.
\]

Dus:
\[
\boxed{\frac{dV}{dt}(5)=25{,}6\pi\ \mathrm{cm^3/s}.}
\]

De kettingregel is nodig omdat:
\[
V\longrightarrow r\longrightarrow t.
\]

\end{oplossing}

\end{oefening}


\begin{oefening}[Temperatuur langs een buis]{\oefU\ESS}

Langs een rechte buis varieert de temperatuur volgens
\[
T(x)=20+5x^2,
\]
waarbij $T$ in graden Celsius en $x$ in meter wordt uitgedrukt.

Een sensor beweegt door de buis volgens
\[
x(t)=0{,}3t.
\]

\begin{question}
Schrijf de temperatuur die de sensor meet als functie van de tijd.
\end{question}

\begin{question}
Bepaal $\dfrac{dT}{dt}$.
\end{question}

\begin{question}
Hoe snel verandert de gemeten temperatuur op $t=10$?
\end{question}

\begin{oplossing}

\[
T(t)=20+5(0{,}3t)^2.
\]

Met de kettingregel:
\[
\frac{dT}{dt}
=
\frac{dT}{dx}\cdot\frac{dx}{dt}.
\]

Nu:
\[
\frac{dT}{dx}=10x
\]
en:
\[
\frac{dx}{dt}=0{,}3.
\]

Dus:
\[
\frac{dT}{dt}=3x.
\]

Voor $t=10$ is:
\[
x=3.
\]

Daarom:
\[
\boxed{\frac{dT}{dt}=9\ ^\circ\mathrm{C/s}.}
\]

\end{oplossing}

\end{oefening}


\begin{oefening}[Een bewegend punt op een parabool]{\oefU\ESS}

Een punt beweegt langs de parabool
\[
y=x^2+1.
\]

Zijn $x$-coördinaat verandert volgens
\[
x(t)=t^2-1.
\]

\begin{question}
Schrijf $y$ als functie van $t$.
\end{question}

\begin{question}
Bepaal $\dfrac{dy}{dt}$ met behulp van de kettingregel.
\end{question}

\begin{question}
Bereken $\dfrac{dy}{dt}$ op het tijdstip $t=2$.
\end{question}

\begin{oplossing}

\[
y(t)=(t^2-1)^2+1.
\]

Met de kettingregel:
\[
\frac{dy}{dt}
=
\frac{dy}{dx}\cdot\frac{dx}{dt}.
\]

Nu:
\[
\frac{dy}{dx}=2x
\]
en:
\[
\frac{dx}{dt}=2t.
\]

Dus:
\[
\frac{dy}{dt}=4tx.
\]

Omdat:
\[
x=t^2-1,
\]
krijgen we:
\[
\boxed{\frac{dy}{dt}=4t(t^2-1).}
\]

Voor $t=2$:
\[
\boxed{\frac{dy}{dt}(2)=24.}
\]

\end{oplossing}

\end{oefening}




\begin{oefening}[Een onbekende parameter]{\oefG}

Gegeven is:
\[
f(x)=(ax+b)^5.
\]

Verder weet je dat:
\[
f'(1)=160
\]
en:
\[
a+b=2.
\]

Bepaal mogelijke waarden van $a$ en $b$.

\begin{oplossing}

\[
f'(x)=5a(ax+b)^4.
\]

Dus:
\[
160=5a(a+b)^4.
\]

Met $a+b=2$:
\[
160=5a\cdot16=80a.
\]

Dus:
\[
a=2
\]
en:
\[
b=0.
\]

Daarom:
\[
\boxed{a=2,\qquad b=0.}
\]

\end{oplossing}

\end{oefening}


\begin{oefening}[Werk achteruit]{\oefG}

Gegeven is:
\[
f'(x)=12x(x^2+1)^5.
\]

Zoek een mogelijke functie $f$ van de vorm
\[
f(x)=A(x^2+1)^n+C.
\]

Bepaal $A$ en $n$.

\begin{oplossing}

\[
f'(x)=2Anx(x^2+1)^{n-1}.
\]

Vergelijk met:
\[
12x(x^2+1)^5.
\]

Daaruit volgt:
\[
n-1=5,
\]
dus:
\[
n=6.
\]

Verder:
\[
2An=12.
\]

Dus:
\[
12A=12,
\]
zodat:
\[
A=1.
\]

Een mogelijke functie is:
\[
\boxed{f(x)=(x^2+1)^6+C.}
\]

\end{oplossing}

\end{oefening}



\begin{oefening}[Zonder volledig uit te werken]{\oefG}

Gegeven:
\[
f(x)=(x^2-4x+5)^6.
\]

\begin{question}
Bepaal $f'(x)$ zonder de zesde macht uit te werken.
\end{question}

\begin{question}
Bepaal alle waarden van $x$ waarvoor
\[
f'(x)=0.
\]
\end{question}

\begin{question}
Leg uit waarom het in deze oefening bijzonder onpraktisch zou zijn
om eerst $f(x)$ volledig uit te werken.
\end{question}

\begin{oplossing}

\[
f'(x)=6(x^2-4x+5)^5(2x-4).
\]

Dus:
\[
\boxed{f'(x)=12(x-2)(x^2-4x+5)^5.}
\]

Voor:
\[
f'(x)=0
\]
geldt:
\[
12(x-2)(x^2-4x+5)^5=0.
\]

Nu:
\[
x^2-4x+5=(x-2)^2+1>0.
\]

Dus:
\[
\boxed{x=2.}
\]

Volledig uitwerken zou een veelterm van graad $12$ opleveren en verbergt
de samengestelde structuur van de functie.

\end{oplossing}

\end{oefening}


% ============================================================
% GEVORDERD — COMBINATIE VAN REGELS
% ============================================================

\begin{oefening}[Product, quotiënt en kettingregel]{\oefG}

Bepaal de afgeleide van:
\[
f(x)=\frac{x^2(x^2+1)^3}{(2x-1)^4}.
\]

Werk systematisch en probeer het eindresultaat te factoriseren.

\begin{oplossing}

We schrijven:
\[
f(x)=x^2(x^2+1)^3(2x-1)^{-4}.
\]

Dan:
\[
\begin{aligned}
f'(x)
&=
2x(x^2+1)^3(2x-1)^{-4}\\
&\quad+
x^2\cdot6x(x^2+1)^2(2x-1)^{-4}\\
&\quad-
8x^2(x^2+1)^3(2x-1)^{-5}.
\end{aligned}
\]

We halen de gemeenschappelijke factor buiten:
\[
f'(x)
=
2x(x^2+1)^2(2x-1)^{-5}
\left[
(x^2+1)(2x-1)+3x^2(2x-1)-4x(x^2+1)
\right].
\]

De uitdrukking tussen haakjes wordt:
\[
4x^3-4x^2+2x-1.
\]

Dus:
\[
\boxed{
f'(x)
=
\frac{
2x(x^2+1)^2(4x^3-4x^2+2x-1)
}
{(2x-1)^5}
}.
\]

\end{oplossing}

\end{oefening}


% ============================================================
% GEVORDERDE VRAAGSTUKKEN
% ============================================================

\begin{oefening}[Een bol smelt]{\oefG}

Een bolvormig stuk ijs heeft op tijdstip $t$ een straal
\[
r(t)=\sqrt{25-2t},
\]
waarbij $r$ in centimeter en $t$ in minuten wordt uitgedrukt.

Het volume van de bol is:
\[
V=\frac{4}{3}\pi r^3.
\]

\begin{question}
Voor welke waarden van $t$ is dit model fysisch betekenisvol?
\end{question}

\begin{question}
Bepaal $\dfrac{dV}{dt}$.
\end{question}

\begin{question}
Leg uit waarom $\dfrac{dV}{dt}$ negatief is.
\end{question}

\begin{question}
Bereken hoe snel het volume afneemt op het tijdstip $t=8$.
\end{question}

\begin{oplossing}

Voor het model moet:
\[
25-2t\geq0.
\]

Dus:
\[
\boxed{0\leq t\leq12{,}5.}
\]

Verder:
\[
V(t)=\frac{4}{3}\pi(25-2t)^{3/2}.
\]

Met de kettingregel:
\[
\frac{dV}{dt}
=
\frac{4}{3}\pi\cdot\frac32(25-2t)^{1/2}\cdot(-2).
\]

Dus:
\[
\boxed{\frac{dV}{dt}=-4\pi\sqrt{25-2t}.}
\]

Voor $t=8$:
\[
25-16=9.
\]

Dus:
\[
\boxed{\frac{dV}{dt}(8)=-12\pi\ \mathrm{cm^3/min}.}
\]

\end{oplossing}

\end{oefening}


\begin{oefening}[Een kegel wordt gevuld]{\oefG}

Een kegelvormig vat heeft hoogte $12\,\mathrm{cm}$ en straal
$6\,\mathrm{cm}$.

Tijdens het vullen vormt het water steeds een gelijkvormige kegel.

Noem $h$ de waterhoogte en $r$ de straal van het wateroppervlak.
Door gelijkvormigheid geldt:
\[
\frac{r}{h}=\frac{6}{12}.
\]

De waterhoogte verandert volgens:
\[
h(t)=\sqrt{t+1}.
\]

Het volume van een kegel is:
\[
V=\frac13\pi r^2h.
\]

\begin{question}
Druk $r$ uit in functie van $h$.
\end{question}

\begin{question}
Schrijf $V$ uitsluitend als functie van $h$.
\end{question}

\begin{question}
Schrijf vervolgens $V$ als functie van $t$.
\end{question}

\begin{question}
Bepaal $\dfrac{dV}{dt}$.
\end{question}

\begin{question}
Welke keten van afhankelijkheden herken je in dit model?
\end{question}

\begin{oplossing}

Uit de gelijkvormigheid:
\[
r=\frac{h}{2}.
\]

Daarom:
\[
V
=
\frac13\pi\left(\frac h2\right)^2h
=
\frac{\pi}{12}h^3.
\]

Omdat:
\[
h(t)=\sqrt{t+1},
\]
krijgen we:
\[
V(t)=\frac{\pi}{12}(t+1)^{3/2}.
\]

Met de kettingregel:
\[
\frac{dV}{dt}
=
\frac{\pi}{12}\cdot\frac32(t+1)^{1/2}.
\]

Dus:
\[
\boxed{\frac{dV}{dt}=\frac{\pi}{8}\sqrt{t+1}.}
\]

De afhankelijkheidsketen is:
\[
t\longrightarrow h\longrightarrow r\longrightarrow V.
\]

Na eliminatie van $r$:
\[
\boxed{t\longrightarrow h\longrightarrow V.}
\]

\end{oplossing}

\end{oefening}


\begin{oefening}[Een camera volgt een drone]{\oefG}

Een drone vliegt horizontaal op een constante hoogte van
$60\,\mathrm{m}$.

De horizontale afstand tot een waarnemer wordt beschreven door:
\[
x(t)=5t,
\]
waarbij $x$ in meter en $t$ in seconden wordt uitgedrukt.

De afstand $d$ tussen de waarnemer en de drone voldoet aan:
\[
d=\sqrt{x^2+60^2}.
\]

\begin{question}
Schrijf $d$ als functie van $t$.
\end{question}

\begin{question}
Bepaal $\dfrac{dd}{dt}$.
\end{question}

\begin{question}
Bereken hoe snel de afstand tot de drone verandert op $t=16$.
\end{question}

\begin{question}
Waarom is $\dfrac{dd}{dt}$ niet gelijk aan $5\,\mathrm{m/s}$?
\end{question}

\begin{oplossing}

\[
d(t)=\sqrt{25t^2+3600}.
\]

Met de kettingregel:
\[
\frac{dd}{dt}
=
\frac{1}{2\sqrt{25t^2+3600}}\cdot50t.
\]

Dus:
\[
\boxed{\frac{dd}{dt}=\frac{25t}{\sqrt{25t^2+3600}}.}
\]

Voor $t=16$:
\[
d(16)=100.
\]

Dus:
\[
\boxed{\frac{dd}{dt}(16)=4\,\mathrm{m/s}.}
\]

De horizontale afstand groeit met $5\,\mathrm{m/s}$, maar $d$ is de schuine
afstand tot de waarnemer. Die verandert dus niet met dezelfde snelheid.

\end{oplossing}

\end{oefening}


\begin{oefening}[Van positie naar vermogen]{\oefG}

De potentiële energie van een systeem wordt in een vereenvoudigd model
gegeven door:
\[
E(x)=3(x^2+4)^2,
\]
waarbij $E$ in joule wordt uitgedrukt.

Een deeltje beweegt volgens:
\[
x(t)=2t-1.
\]

Het vermogen dat samenhangt met de verandering van deze energie wordt
gedefinieerd als:
\[
P(t)=\frac{dE}{dt}.
\]

\begin{question}
Schrijf $E$ als functie van $t$.
\end{question}

\begin{question}
Bepaal $P(t)$.
\end{question}

\begin{question}
Bereken het vermogen op $t=2$.
\end{question}

\begin{question}
Schrijf de berekening ook in de vorm
\[
\frac{dE}{dt}
=
\frac{dE}{dx}\cdot\frac{dx}{dt}.
\]
Leg uit wat beide factoren fysisch voorstellen.
\end{question}

\begin{oplossing}

\[
\frac{dE}{dx}
=
3\cdot2(x^2+4)\cdot2x
=
12x(x^2+4).
\]

Verder:
\[
\frac{dx}{dt}=2.
\]

Daarom:
\[
P(t)=24x(x^2+4).
\]

Met:
\[
x=2t-1
\]
volgt:
\[
\boxed{
P(t)
=
24(2t-1)\left((2t-1)^2+4\right)
}.
\]

Voor $t=2$:
\[
x=3.
\]

Dus:
\[
\boxed{P(2)=936\,\mathrm W.}
\]

De factor $\frac{dE}{dx}$ beschrijft hoe de energie verandert met de positie.
De factor $\frac{dx}{dt}$ is de snelheid waarmee de positie verandert.

\end{oplossing}

\end{oefening}


\begin{oefening}[Omgekeerd redeneren bij een groeiende cirkel]{\oefG}

De oppervlakte van een cirkel is:
\[
A=\pi r^2.
\]

Op een bepaald ogenblik is:
\[
r=5\,\mathrm{cm}.
\]

Op dat ogenblik groeit de oppervlakte met:
\[
20\pi\,\mathrm{cm^2/s}.
\]

\begin{question}
Bepaal op dat ogenblik $\dfrac{dr}{dt}$.
\end{question}

\begin{question}
Leg uit waarom dit een toepassing van de kettingregel is, hoewel
$r(t)$ niet expliciet gegeven is.
\end{question}

\begin{oplossing}

Met de kettingregel:
\[
\frac{dA}{dt}
=
\frac{dA}{dr}\cdot\frac{dr}{dt}.
\]

Omdat:
\[
A=\pi r^2,
\]
geldt:
\[
\frac{dA}{dr}=2\pi r.
\]

Dus:
\[
20\pi
=
2\pi\cdot5\cdot\frac{dr}{dt}.
\]

Daarom:
\[
\boxed{\frac{dr}{dt}=2\,\mathrm{cm/s}.}
\]

Hoewel $r(t)$ niet expliciet gekend is, hangt $A$ via $r$ af van de tijd:
\[
A=A(r(t)).
\]

\end{oplossing}

\end{oefening}


\end{document}